% Part: first-order-logic
% Chapter: semantics
% Section: intro-semantics

\documentclass[../../../include/open-logic-section]{subfiles}

\begin{document}

\olfileid{fol}{syn}{its}

\olsection{परिचय}

व्यंजकांना अर्थ देणे हे चिन्हार्थमीमांसेचे क्षेत्र आहे. चिन्हार्थमीमांसेतील
मध्यवर्ती संकल्पना म्हणजे !!a{structure} मधील पूर्ति. !!^a{structure}
भाषेच्या मूलभूत घटकांना अर्थ देते: !!a{domain} म्हणजे वस्तूंचा अरिक्त
संच. संख्यापक या प्रांतातील घटकांवर फिरतात असा त्यांचा अर्थ लावला जातो,
!!{constant}s ना प्रांतातील घटक नेमून दिले जातात, !!{function}s ना
!!{domain} पासून त्याच प्रांताकडे जाणारी फलने नेमून दिली जातात आणि
!!{predicate}s ना !!{domain} वरील संबंध नेमून दिले जातात. !!{domain}
आणि मूलभूत शब्दसंग्रहातील चिन्हांना केलेल्या नेमणुका मिळून
!!a{structure} बनते.
!!^{variable}s ही !!{formula}s मध्ये येऊ शकतात; त्यांना चिन्हार्थ देण्यासाठी
आपल्याला !!{domain} मधील !!{element}s ही नेमून द्यावी लागतात---हेच
चर-मूल्यांकन होय. शेवटी, पूर्ति-संबंध या सर्व गोष्टी एकत्र आणतो.
!!^a{formula} ची !!a{structure}~$\Struct{M}$ मध्ये !!a{variable} च्या
मूल्यांकन~$s$ च्या सापेक्ष पूर्ति होऊ शकते; हे $\Sat{M}{!A}[s]$ असे
लिहितात. हा संबंधदेखील~$!A$ च्या संरचनेवरील विगमनाने परिभाषित केला जातो.
उदाहरणार्थ, $(!A \land !B)$ ची पूर्ति परिभाषित करताना तार्किक संयोजकांचे
सत्यता-कोष्टक वापरून ती $!A$ आणि~$!B$ यांची पूर्ति होते (किंवा होत नाही)
यांच्या आधारे सांगितली जाते. त्यानंतर असे निष्पन्न होते की !!{formula}~$!A$
हे !!a{sentence}, म्हणजे त्यात मुक्त चरे नसतील, तर !!{variable} चे
मूल्यांकन असंबद्ध ठरते; त्यामुळे !!{structure}s मध्ये !!{sentence}s ची
सरळ पूर्ति होते (किंवा होत नाही) असे आपण म्हणू शकतो.

!!{sentence}s साठीच्या पूर्ति-संबंध~$\Sat{M}{!A}$ च्या आधारे आपण वैधता,
तार्किक निष्पन्नता आणि पूर्ततायोग्यता या मूलभूत चिन्हार्थविषयक संकल्पना
परिभाषित करू शकतो. !!^a{sentence} ची प्रत्येक !!{structure} मध्ये पूर्ति
होत असेल, तर ते वैध, $\Entails !A$, असते. !!{sentence}s च्या एखाद्या
संचातून ते तार्किकरीत्या निष्पन्न होते, $\Gamma \Entails !A$, जर~$\Gamma$ मधील
सर्व !!{sentence}s ची पूर्ति करणारे प्रत्येक !!{structure} हे~$!A$ चीही
पूर्ति करत असेल. आणि एखाद्या !!{sentence}s च्या संचातील सर्व
!!{sentence}s ची एकाच वेळी पूर्ति करणारे एखादे !!{structure} असेल, तर तो
संच पूर्ततायोग्य असतो. !!{formula}s विगमनाधारित रीतीने परिभाषित केली आहेत
आणि पूर्ति हीही !!{formula}s च्या संरचनेवरील विगमनाने परिभाषित केली जाते;
म्हणून आपल्या चिन्हार्थमीमांसेचे गुणधर्म सिद्ध करण्यासाठी आणि परिभाषित
चिन्हार्थविषयक संकल्पनांचा परस्परसंबंध लावण्यासाठी विगमन वापरता येते.

\end{document}
